% Generated draft: requires mathematical and source review.
\documentclass[11pt]{article}
\usepackage[margin=25mm]{geometry}
\usepackage{fontspec}
\setmainfont{Latin Modern Roman}
\usepackage{amsmath,amssymb,mathtools,unicode-math}
\setmathfont{Latin Modern Math}
\usepackage{enumitem,needspace,xurl,hyperref}
\setlength{\parindent}{0pt}
\setlength{\parskip}{5pt}
\setlength{\emergencystretch}{3em}
\newcommand{\N}{\mathbb{N}}
\newcommand{\Z}{\mathbb{Z}}
\newcommand{\Q}{\mathbb{Q}}
\newcommand{\R}{\mathbb{R}}
\newcommand{\C}{\mathbb{C}}
\newcommand{\F}{\mathbb{F}}
\newcommand{\A}{\mathbb{A}}
\newcommand{\PP}{\mathbb{P}}
\newcommand{\E}{\mathbb{E}}
\newcommand{\eps}{\varepsilon}
\newcommand{\dd}{\,\mathrm{d}}
\begin{document}
\section*{UnsolvedMath ID: 20002719}
\textbf{Source problem code:} catalogue entry
\subsection{Definitions and mathematical statement}
F. Goodman and P. Sniady's Problem 5.1 asks for a natural model of the type-B free probability introduced by Philippe Biane, Fred Goodman, and Alexandru Nica, and for analogues in other reflection types. Its historical assertion that no model was then known is not a current status claim. The accompanying remark stresses that usual free-probability formulas use the whole sequence \(NC(A_{n-1})\), mentions Reiner's type-B multiplicative functions, and asks how an exceptional type without an infinite family could support the required all-arity structure.

The aid reports later identification of scalar first-order type-B content with infinitesimal free probability and finite-rank random-matrix perturbation models, without complete citations. An infinitesimal law here is \((\varphi,\varphi')\) with \(\varphi\) a unital state and \(\varphi'(1)=0\), arising as \(\varphi_N=\varphi+N^{-1}\varphi'+o(N^{-1})\). The derivative \(\varphi'\) need not itself be positive. Infinitesimal cumulants are derivatives of free cumulants; infinitesimal freeness means ordinary and first-order mixed free cumulants vanish. This specifies the scalar first-order interpretation, not every possible type-B bimodule construction.

The aid's explicit zero-bulk model fixes \(r\) mutually orthogonal rank-one projections \(P_{s,N}\in M_N(\C)\), \(N\ge r\), real spike values \(\theta_s\), and a partition of the spike labels into coordinate sets \(I_j\). Put
\[
 A_{j,N}=\sum_{s\in I_j}\theta_sP_{s,N},\qquad\tau_N=N^{-1}\operatorname{Tr}.
\]
For every noncommutative polynomial \(P\), the exact identity is
\[
 \tau_NP(A_{1,N},\ldots,A_{k,N})
 =P(0)+\frac1N\sum_j\sum_{s\in I_j}
 \bigl(P(0,\ldots,\theta_s,\ldots,0)-P(0)\bigr),
\]
where the nonzero scalar slot is \(j\). The families have ordinary zero laws; mixed words using different coordinates have zero first-order moments and cumulants. Thus they are infinitesimally free in this special model, although their finite matrices commute.

For one combined spike variable with all \(r\) spikes, the empirical law and infinitesimal law are
\[
 \mu_N=(1-r/N)\delta_0+N^{-1}\sum_{s=1}^r\delta_{\theta_s},\quad
 \varphi=\delta_0,\quad\varphi'=\sum_s\delta_{\theta_s}-r\delta_0,
\]
and for every \(n\ge1\),
\[
 \kappa_n'=\sum_{s=1}^r\theta_s^n.
\]
If \(r\) is known, the first \(r\) such cumulants determine the multiset of all spikes by Newton identities. Explicitly, for \(p_j=\kappa_j'\), \(e_0=1\),
\[
 ke_k=\sum_{j=1}^k(-1)^{j-1}e_{k-j}p_j,
\]
and the spikes are the roots, with multiplicity, of \(z^r-e_1z^{r-1}+\cdots+(-1)^re_r\). The known-count hypothesis includes any zero spikes, which positive-order power sums alone cannot count if \(r\) is unknown. Every finite-\(N\) trace is a positive tracial state even though its first-order correction is signed. This special deterministic construction is not a universal random-matrix realization of all type-B laws, nor an answer for every exceptional type; the aid reports that broader all-arity issue unresolved.
\subsection{Short English statement}
The aid distinguishes the historical type-B modelling request from later infinitesimal interpretations. Its orthogonal-spike example gives an exact positive finite-matrix expansion, infinitesimal freeness, and power-sum cumulants recovering a known number of spikes. It does not resolve the all-arity exceptional-type question.
\subsection{Sources}
\begin{itemize}
\item Local registry, permanent UnsolvedMath ID \(20002719\), exact public code \texttt{catalogue entry}; canonical location \texttt{aim-probability-notes.json}, \texttt{notes[160]}.
\item F. Goodman and P. Sniady, Problem 5.1, AIM \emph{Braid groups, clusters and free probability}; supplied reference \url{https://aimath.org/WWN/braidgroups/braidgroups.pdf}. The verification field reports prior inspection of page 10 correcting the split number and Alexandru Nica's name. No PDF was accessed here.
\item Source references [10] for Biane--Goodman--Nica type B, [48] for multiplicative-function formulas, and [42] for Vic Reiner's type-B multiplicative functions are preserved as unresolved reference labels; their full entries are not supplied. Later infinitesimal and finite-rank random-matrix work is reported without bibliographic details.
\item Supplied research attempt 1: \texttt{partial\_result}, run and registry status \texttt{partially\_solved}, rigor label \texttt{proved}, low novelty confidence; classification \texttt{PARTIAL-PROGRESS}, created and updated 14 August 2026. The aid is machine generated, not peer reviewed, and requires independent verification before citation.
\end{itemize}
\end{document}
